\subsection*{10.4 Convergence of Random Vectors}

In this section, we formally establish the relationship between the convergence of random vectors and the convergence of their component random variables.

Consider measurable functions that take value in $\mathbb{R}^d$. Let $(\Omega,\mathcal{F},P)$ denote a probability space and $V(\omega)$ denote a measurable mapping from $(\Omega,\mathcal{F})$ to $(\mathbb{R}^d,\mathcal{B}(\mathbb{R}^d))$. We can prove that a vector function is measurable if and only if each component is a measurable function. This result in measure theory is known as the coordinate-wise convergence theorem.

\begin{thmbox}{10.9}
Let $(\Omega,\mathcal{F})$ be a measurable space and
\[
V(\omega)\coloneqq (X_1(\omega),X_2(\omega),\ldots,X_d(\omega))
\]
be a mapping from $\Omega$ to $\mathbb{R}^d$, where $X_i(\omega)$ is the $i$-th component function. Then $V(\omega)$ is $(\mathcal{F},\mathcal{B}(\mathbb{R}^d))$-measurable if and only if $X_i(\omega)$ is $(\mathcal{F},\mathcal{B}(\mathbb{R}))$-measurable for $i=1,2,\ldots,d$.
\end{thmbox}

\textit{Proof}
Suppose $X_i$ is measurable for all $i$. We want to show that $V^{-1}(B)$ is $\mathcal{F}$-measurable for any Borel set $B$ in $\mathcal{B}(\mathbb{R}^d)$. Since $\mathcal{B}(\mathbb{R}^d)$ is generated by $d$-dimensional hyper-cubes (see Theorem 2.8), it is sufficient to show that $V^{-1}(C)$ is $\mathcal{F}$-measurable for any hyper-cube $C$ of the form
\[
(a_1,b_1)\times(a_2,b_2)\times\cdots\times(a_d,b_d),
\]
where $a_i,b_i$ are real numbers, possibly $\pm\infty$.

For any hyper-cube $C$, the inverse image $V^{-1}(C)$ is equal to
\[
\{\omega:X_i(\omega)\in(a_i,b_i)\text{ for }i=1,2,\ldots,d\}
=
\bigcap_{i=1}^{d}\{\omega:X_i(\omega)\in(a_i,b_i)\}.
\]
Since each set in the intersection is $\mathcal{F}$-measurable, the set $V^{-1}(C)$ is also $\mathcal{F}$-measurable.

Conversely, suppose $V$ is $\mathcal{F}$-measurable. That is, we suppose that $V^{-1}(B)$ is $\mathcal{F}$-measurable for all Borel sets $B\in\mathcal{B}(\mathbb{R}^d)$. For each $i=1,2,\ldots,d$, let $A_i$ be the cylinder
\[
A_i=\mathbb{R}\times\cdots\times\mathbb{R}\times(a_i,b_i)\times\mathbb{R}\times\cdots\times\mathbb{R},
\]
where the interval $(a_i,b_i)$ is in the $i$-th component. This is a Borel set. Then $V^{-1}(A_i)$ is $\mathcal{F}$-measurable, and it is equal to
\[
V^{-1}(A_i)=X_i^{-1}((a_i,b_i)).
\]
Hence, $X_i^{-1}((a_i,b_i))$ is $\mathcal{F}$-measurable for all $a_i,b_i\in\mathbb{R}$.
\hfill $\square$

We measure distance in $\mathbb{R}^d$ by the Euclidean norm and use the following notation:
\[
\lVert(x_1,x_2,\ldots,x_d)\rVert
\coloneqq
\sqrt{x_1^2+x_2^2+\cdots+x_d^2}.
\]

\begin{defbox}{10.6}
Let $(\Omega,\mathcal{F},P)$ denote a probability space, and let $V_n(\omega)$ be a sequence of $d$-dimensional random vectors defined on $\Omega$. The sequence of random vectors $V_n(\omega)$ is said to converge almost surely to $V(\omega)$ if there is an event $E$ with $P(E)=1$ such that
\[
\lim_{n\to\infty}V_n(\omega)=V(\omega)
\]
for all $\omega\in E$.

We say that $V_n(\omega)$ converges in probability to $V(\omega)$ if for all $\epsilon>0$, the probability
\[
P(\{\omega:\lVert V_n(\omega)-V(\omega)\rVert>\epsilon\})
\]
approaches zero as $n\to\infty$.
\end{defbox}

The next theorem gives the relation between convergence as a random vector and convergence in each individual component.

\begin{thmbox}{10.10}
Let $(\Omega,\mathcal{F},P)$ be a probability space. For $n=1,2,3,\ldots$, let
\[
V_n(\omega)=(X_{n1}(\omega),X_{n2}(\omega),\ldots,X_{nd}(\omega))
\]
be a $d$-dimensional random vector, and let
\[
V(\omega)=(X_1(\omega),X_2(\omega),\ldots,X_d(\omega))
\]
be another $d$-dimensional random vector. Then:
\begin{enumerate}
\item $V_n\xrightarrow{\mathrm{a.s.}}V$ if and only if $X_{ni}\xrightarrow{\mathrm{a.s.}}X_i$ for all $i=1,\ldots,d$.
\item $V_n\xrightarrow{P}V$ if and only if $X_{ni}\xrightarrow{P}X_i$ for all $i=1,\ldots,d$.
\end{enumerate}
\end{thmbox}

\textit{Proof}
For part (a), suppose the component functions converge a.s. That is, for each $i=1,2,\ldots,d$, we can find an event $E_i$ with $P(E_i)=1$ such that
\[
X_{ni}(\omega)\xrightarrow{\mathrm{a.s.}}X_i(\omega)
\]
for $\omega\in E_i$. Then in $\bigcap_{i=1}^{d}E_i$, we have the convergence of random vector,
\[
\lim_{n\to\infty}V_n(\omega)=V(\omega)
\qquad\text{for }\omega\in\bigcap_{i=1}^{d}E_i.
\]
The proof is finished by noting that $P(\bigcap_{i=1}^{d}E_i)=1$.

Conversely, suppose $V_n\xrightarrow{\mathrm{a.s.}}V$. There is a set $E$ with probability $1$ such that $V_n(\omega)\to V(\omega)$ for $\omega\in E$. Each component of $V_n(\omega)$ converges to the corresponding component of $V$ for $\omega\in E$.

For part (b), assume that each component of $V_n$ converges in probability to the corresponding component of $V$. For any $\omega\in\Omega$,
\[
\lvert X_{ni}(\omega)-X_i(\omega)\rvert\leq \frac{\epsilon}{\sqrt{d}}
\text{ for all } i
\quad\Longrightarrow\quad
\lVert V_n(\omega)-V(\omega)\rVert\leq \epsilon.
\]
We thus have the following set inclusion:
\[
\{\omega:\lVert V_n-V\rVert>\epsilon\}
\subseteq
\bigcup_{i=1}^{d}
\left\{\omega:\lvert X_{ni}-X_i\rvert>\frac{\epsilon}{\sqrt{d}}\right\}.
\]
Taking probability of both sides, we get $P(\lVert V_n-V\rVert>\epsilon)\to 0$ as $n\to\infty$, which follows from the union bound (Exercise 2.4).

Conversely, assume that $V_n$ converges to $V$ in probability as a sequence of random vectors. Consider the $i$-th component of $V_n$. Since
\[
\lvert X_{ni}(\omega)-X_i(\omega)\rvert
\leq
\lVert V_n(\omega)-V(\omega)\rVert,
\]
we have
\[
\{\omega:\lvert X_{ni}(\omega)-X_i(\omega)\rvert>\epsilon\}
\subseteq
\{\omega:\lVert V_n(\omega)-V(\omega)\rVert>\epsilon\}.
\]
Hence, for any $\epsilon>0$,
\[
V_n\xrightarrow{P}V
\quad\Longrightarrow\quad
P(\lVert V_n(\omega)-V(\omega)\rVert>\epsilon)\to 0
\]
as $n\to\infty$, which implies
\[
P(\lvert X_{ni}(\omega)-X_i(\omega)\rvert>\epsilon)\to 0
\]
as $n\to\infty$. By the definition of convergence in probability, $X_{ni}$ converges to $X_i$ in probability for all $i$.
\hfill $\square$

Using Theorem 10.10, we can check whether a sequence of random vectors converges in probability by considering the convergence of each component.
